\documentclass[12pt]{article}
\usepackage[utf8]{inputenc}
\usepackage[T1]{fontenc}
\usepackage[margin=1in]{geometry}
\usepackage{setspace}
\usepackage{xcolor}
\usepackage[colorlinks=true, urlcolor=blue!60!black, linkcolor=blue!60!black]{hyperref}
\setlength{\parindent}{2em}
\setlength{\parskip}{1.25em}
\renewcommand{\baselinestretch}{1.0}
\begin{document}
\begin{center}
{\LARGE \textit{Trump v. United States}, 603 U.S. 593 (2024)\par}
\vspace{0.8em}
{\large Case Brief --- Four-Page Edition\par}
\end{center}
\vspace{1em}

\section{Memory Jogger}

A former President enjoys absolute criminal immunity for exercises of core, exclusive constitutional powers, at least presumptive immunity for all other official acts within the outer perimeter of presidential responsibility, and no immunity for unofficial acts (\href{https://www.law.cornell.edu/supremecourt/text/23-939}{603 U.S.\ 593}).

\section{Detailed Case Facts}

After Donald J.\ Trump lost the 2020 election, Special Counsel Jack Smith, appointed in November 2022, obtained a four-count indictment on August 1, 2023, alleging that Trump conspired to overturn the election through knowingly false claims of fraud. The alleged means were pressuring state officials to change vote counts and replace electors; organizing fraudulent slates of electors in seven states; trying to use the Justice Department to conduct sham investigations and send a false letter to states, while threatening to replace the Acting Attorney General; pressuring Vice President Pence to reject electoral votes on January 6, 2021, and directing supporters to the Capitol; and exploiting the Capitol attack to delay certification. The charges were conspiracy to defraud the United States, conspiracy to obstruct and obstruction of an official proceeding, and conspiracy against rights.

\section{Procedural History}

Judge Tanya Chutkan denied Trump's motion to dismiss on immunity and Impeachment Judgment Clause grounds. The Court denied certiorari before judgment, and a unanimous D.C.\ Circuit panel affirmed, reasoning that ``former President Trump has become citizen Trump.'' The Supreme Court granted review, heard argument on April 25, 2024, and on July 1, 2024, vacated and remanded. After Trump's reelection, the Justice Department moved to dismiss, and the case was dismissed on November 25, 2024.

\section{Judicial Votes}

Vacated and remanded, \textbf{6--3}; \textbf{5--4} as to Part III-C.
\begin{itemize}
\item \textbf{Majority:} Chief Justice Roberts, joined by Justices Thomas, Alito, Gorsuch, and Kavanaugh, and by Justice Barrett except Part III-C.
\item \textbf{Concurring:} Justice Thomas; Justice Barrett, concurring in part.
\item \textbf{Dissenting:} Justice Sotomayor, joined by Justices Kagan and Jackson; Justice Jackson.
\end{itemize}

\section{Holding}

A former President is absolutely immune for conduct within his ``conclusive and preclusive'' constitutional authority, at least presumptively immune for other official acts within the ``outer perimeter'' of his responsibility (rebuttable only if prosecution poses no ``dangers of intrusion on the authority and functions of the Executive Branch''), and not immune for unofficial acts. Courts may not inquire into motive or deem an act unofficial merely because it allegedly violates a criminal law. Part III-C barred evidence of immune official acts even in prosecutions for unofficial conduct. Trump was absolutely immune for his Justice Department communications and presumptively immune for his dealings with Pence; the classification of the state-official, elector, and January 6 allegations was remanded. The Court rejected the claim that impeachment and conviction must come first and reserved whether non-core immunity is absolute.

\section{Analysis of Opinions}

\textbf{Majority (Roberts).} Drawing on Justice Jackson's \textit{Youngstown} framework, Roberts held that Congress cannot act on subjects within the President's exclusive powers, such as pardons, recognition, removal, and directing prosecutions, so courts cannot adjudicate prosecutions examining them. For other official acts, Hamilton's call for an ``energetic'' Executive, \textit{Nixon v.\ Fitzgerald}, and the executive-privilege cases supported presumptive immunity, because the threat of prosecution would chill ``bold and unhesitating action'' and risk ``an Executive Branch that cannibalizes itself.'' Allowing juries to hear evidence of immune acts would let them judge that conduct indirectly.

\textbf{Thomas (concurring).} Thomas questioned whether the Special Counsel's office was ``established by Law'' under the Appointments Clause, a view that influenced the later dismissal of the classified-documents case.

\textbf{Barrett (concurring in part).} Immunity means an as-applied constitutional challenge plus interlocutory review. She would apply ordinary statutory interpretation and then ask whether prosecution intrudes on executive authority, viewed organizing fraudulent electors as private conduct, and rejected the evidentiary bar, which would cripple bribery prosecutions.

\textbf{Sotomayor (dissenting).} The immunity is ``atextual, ahistorical, and unjustifiable'': the Framers granted legislative but not presidential immunity, the Impeachment Judgment Clause contemplates later prosecution, and Nixon's acceptance of a pardon confirms the understanding. A President who orders SEAL Team 6 to kill a rival would be ``[i]mmune, immune, immune.'' ``With fear for our democracy, I dissent.''

\textbf{Jackson (dissenting).} The Court replaced an ``individual accountability model'' with a ``Presidential accountability model,'' reviving the notion that ``the King can do no wrong'' and making the Court the gatekeeper of presidential prosecutions.

\section{Examples: Future Applications}

\begin{enumerate}
\item \textbf{Same-Side (Immune): Corruptly Halting an Investigation.} Directing the Attorney General to drop a donor's fraud case, on threat of firing, falls within conclusive and preclusive authority, and motive is off limits.
\item \textbf{Same-Side (Immune): Pardoning a Witness.} Pardoning a Cabinet official to prevent testimony is an exercise of the pardon power Congress cannot criminalize.
\item \textbf{Opposite-Side (Not Immune): Business Fraud.} Falsifying a family business's tax returns and loan applications is unofficial, as in \textit{Clinton v.\ Jones}.
\item \textbf{Opposite-Side (Not Immune): Private Campaign Wiretap.} Paying a private investigator to wiretap a rival's headquarters is conduct as a candidate, not an Article II act.
\item \textbf{Fence-Sitter: Foreign Aid for an Investigation.} Conditioning appropriated aid on a foreign leader's investigation of a rival involves diplomacy within the outer perimeter, but Congress shares authority over appropriations, and Part III-C would make proof extraordinarily difficult.
\end{enumerate}

\section{Critique}

\textbf{Originalist and textualist critiques.} Jack Goldsmith calls the Court's account of exclusive power over prosecutions and removal among the most far-reaching in its history, resting on functional fears rather than text; William Baude notes the reasoning went beyond any determinate text or tradition. Shalev Gad Roisman argues that the Court labeled powers ``exclusive'' without originalist support and then abandoned formalism for balancing. Jonathan Adler defends a narrower immunity like Barrett's as following from Article II.

\textbf{Progressive and doctrinal critiques.} Gillian Metzger argues that the Court engaged in ``faux balancing,'' sidelined Congress, and weakened internal executive-branch legal guardrails. Critics contrast the Court's strict text-and-history method in \textit{Dobbs} and \textit{Bruen} with its consequentialist structuralism here.

\section{Key Quotations}

\begin{quote}
``We conclude that under our constitutional structure of separated powers, the nature of Presidential power requires that a former President have some immunity from criminal prosecution for official acts during his tenure in office. At least with respect to the President's exercise of his core constitutional powers, this immunity must be absolute. As for his remaining official actions, he is also entitled to immunity.''
\end{quote}
\noindent --- Chief Justice Roberts, opinion of the Court.

\begin{quote}
``The President enjoys no immunity for his unofficial acts, and not everything the President does is official. The President is not above the law. But Congress may not criminalize the President's conduct in carrying out the responsibilities of the Executive Branch under the Constitution.''
\end{quote}
\noindent --- Chief Justice Roberts, opinion of the Court.

\begin{quote}
``The relationship between the President and the people he serves has shifted irrevocably. In every use of official power, the President is now a king above the law.''
\end{quote}
\noindent --- Justice Sotomayor, dissenting.

\end{document}
